% =====================================================================
%  UQFF Yang-Mills Mass Gap Derivation — arxiv preprint (FILLED)
%
%  Target submission categories:  math-ph (primary), hep-th (cross-list)
%
%  Compile with:
%      pdflatex preprint_filled.tex
%      pdflatex preprint_filled.tex     (twice for cross-refs)
%
%  Author:    Daniel T. Murphy
%  License:   AGPL-3.0-or-later (or commercial; see Star-Magic repo)
% =====================================================================

\documentclass[11pt,a4paper,onecolumn]{article}

\usepackage[utf8]{inputenc}
\usepackage[T1]{fontenc}
\usepackage{amsmath, amssymb, amsthm}
\usepackage{graphicx}
\usepackage{hyperref}
\usepackage{url}
\usepackage[margin=1in]{geometry}
\usepackage{microtype}
\usepackage{booktabs}
\usepackage{listings}
\usepackage{xcolor}
\usepackage{authblk}

\hypersetup{
  colorlinks=true,
  linkcolor=blue!50!black,
  citecolor=blue!50!black,
  urlcolor=blue!60!black
}

\lstset{
  basicstyle=\ttfamily\small,
  breaklines=true,
  frame=single,
  framerule=0.2pt,
  rulecolor=\color{gray!30},
  numbers=none,
  language=Python
}

\theoremstyle{plain}
\newtheorem{theorem}{Theorem}
\newtheorem{proposition}[theorem]{Proposition}
\newtheorem{lemma}[theorem]{Lemma}

\theoremstyle{definition}
\newtheorem{definition}[theorem]{Definition}
\newtheorem{claim}[theorem]{Claim}

\theoremstyle{remark}
\newtheorem{remark}[theorem]{Remark}

\title{
  Positive-Definite Energy Threshold from the UQFF Vacuum Ledger \\
  as a Candidate Physical Realization of the Yang--Mills Mass Gap
}

\author[1]{Daniel T. Murphy}
\affil[1]{
  Star-Magic Research Program, Youngstown, OH, USA \\
  \texttt{daniel.murphy00@enrgyone.com}
}

\date{July 2026}

\begin{document}

\maketitle

\begin{abstract}
We derive a positive-definite energy threshold
$E_{\text{crack}} = \rho_{\mathrm{SCm}} c^2 / [\mathrm{SSq}]
                  = 1.118 \times 10^{-19}\,\mathrm{J}$
($\approx 0.7\,\mathrm{eV}$ at the formula level)
from the Unified Quantum Field Framework (UQFF) closed vacuum ledger
using only two locked primitives --- the SuperConductive Material
vacuum energy density $\rho_{\mathrm{SCm}}$ and the canonical resonance
ratio $[\mathrm{SSq}]$ --- together with the speed of light $c$. The
derivation requires zero free parameters and zero fitting. Because
each factor on the right-hand side is positive by physical definition,
$E_{\text{crack}} > 0$ holds by construction. We propose this as a
candidate physical realization of the Yang--Mills mass gap and discuss
its multi-designation cluster-position landscape extending from the
sub-eV vacuum-cracking floor through a $\sim 700\,\mathrm{eV}$
DPM-vortex experimental regime to a $\sim 1.736\,\mathrm{GeV}$
lightest-glueball scale consistent with lattice-QCD numerical
evidence~\cite{douglas2026}. We acknowledge that this constitutes a
physics-level proposal rather than a Clay-compliant
Wightman-axiom-verified mathematical construction, and identify the
operator-algebraic translation as the principal future-work step.
The full framework is reproducibly available via
\verb|pip install uqff==5.35.0|.
\end{abstract}

\bigskip

\noindent
\textbf{Keywords}: Yang--Mills mass gap; vacuum ledger; positive-definite
energy threshold; Di-Pseudo-Monopole vortex; SCm/UA substrate;
multi-designation cluster positions; lattice-QCD consistency; Clay
Millennium Prize.

\medskip

\noindent
\textbf{MSC 2020}: 81T13 (Yang--Mills and other gauge theories), 81T08
(Constructive QFT), 81V25 (Other elementary particle theory).

% =====================================================================
\section{Introduction}\label{sec:intro}
% =====================================================================

\subsection{Problem statement}

The Clay Mathematics Institute's Yang--Mills existence and mass gap
problem~\cite{jaffe2000} requires two tightly linked results:

\begin{enumerate}
\item[(E)] \textbf{Existence.} Construction of a rigorous mathematical
  theory of quantum Yang--Mills fields on four-dimensional spacetime
  for a compact simple gauge group (typically $\mathrm{SU}(2)$ or
  $\mathrm{SU}(3)$), satisfying the Wightman axioms or an equivalent
  system.
\item[(G)] \textbf{Mass gap.} A proof that this theory has a positive
  mass gap $\Delta > 0$: the lightest particle in the spectrum has
  strictly positive mass, with a spectral gap between the vacuum energy
  and the first excited state.
\end{enumerate}

Informally: prove that the gauge bosons (gluons, in the QCD context)
acquire effective mass through non-perturbative effects, even though
the classical Yang--Mills Lagrangian is massless. This mechanism
underpins the short-range character of the strong interaction and the
confinement of quarks into hadrons.

\subsection{Current state of progress (as of mid-2026)}

As of the time of writing, the problem remains officially unsolved and
the \$1{,}000{,}000 Clay prize is unclaimed. Lattice QCD simulations
strongly suggest a mass gap on the order of hundreds of MeV --- with
the lightest $0^{++}$ glueball estimated at approximately $1.6$--$1.7$
GeV in pure $\mathrm{SU}(3)$ Yang--Mills theory --- but this
numerical evidence is not a rigorous mathematical proof.

Significant advances have occurred in lower dimensions. Hairer and
collaborators~\cite{hairer2024} rigorously constructed the
stochastic quantisation of Yang--Mills--Higgs in three dimensions using
regularity structures, following Hairer's earlier 2022 result on Langevin
dynamics for the two-dimensional Yang--Mills measure. Douglas's
January 2026 review~\cite{douglas2026} in
\emph{Nature Reviews Physics} surveys the landscape and highlights
stochastic quantisation as one of the most promising directions.
Complementary approaches under active investigation include
constructive quantum field theory, gauge-invariant regularisation
schemes, and functional-integral methods with rigorous control of the
continuum limit. The four-dimensional case --- which is what Clay
requires --- remains open.

\subsection{This paper}

We propose that the Unified Quantum Field Framework (UQFF) provides a
concrete non-perturbative mechanism generating a positive-definite
energy threshold, which we identify with the Yang--Mills mass gap in
its physical rather than operator-algebraic sense. The mechanism ---
Di-Pseudo-Monopole (DPM) vortex formation at the interface of maximum
attraction between the Universal Aether (UA) and the SuperConductive
Material (SCm) substrates --- yields an energy threshold
$E_{\text{crack}} = \rho_{\mathrm{SCm}} c^2 / [\mathrm{SSq}]$ that is
positive by construction from two locked pre-cosmogenic primitives.

We stress that the present submission is a \emph{physics-level
proposal}, not a Clay-compliant mathematical proof. The
translation from the SCm/UA dynamical substrate to a
Wightman-axiom-compliant operator algebra on four-dimensional
Minkowski space is identified as the principal open mathematical
question and is scoped as future work in \S\ref{sec:future}. Our
contribution is: (i) an explicit closed-form derivation of a
non-fitted, positive-definite mass-gap-scale energy from locked
primitives (\S\ref{sec:derivation}); (ii) a multi-cluster spectral
landscape consistent with lattice-QCD numerical evidence
(\S\ref{sec:clusters}, \S\ref{sec:lattice}); (iii) a public,
reproducible computational framework
(\S\ref{sec:reproducibility}); and (iv) a falsifiable experimental
prediction in the high-field vacuum regime (\S\ref{sec:scope}).

% =====================================================================
\section{The UQFF Vacuum Ledger}\label{sec:ledger}
% =====================================================================

The Unified Quantum Field Framework is a vacuum-first theoretical
structure resting on nine truly-independent primitives: the integer
lattice $(D_{\mathrm{phys}}, D_{\mathrm{crit}}, N_{\mathrm{ch}},
\mathrm{SO}_5, A_5) = (4, 26, 9, 10, 60)$ and four real primitives
$(\rho_{\mathrm{SCm}}, \beta_i, \Phi_{\mathrm{res}}, F_{\mathrm{TRZ}})$.
Two structural quantities used elsewhere in the framework
($D_{\mathrm{BSFG}} = 6$ and $K_{\mathrm{Mex}} = 25/12$) are derivable
from these nine via the relations $D_{\mathrm{BSFG}} =
D_{\mathrm{crit}} - 2\,\mathrm{SO}_5$ and $K_{\mathrm{Mex}} =
\Phi_{5/6} \cdot \mathrm{SO}_5 / D_{\mathrm{phys}}$~\cite{murphy2026paper1521, murphy2026paper1522}, so the truly-independent count is nine.
For the derivation in this paper, only three quantities are required.

\begin{definition}[Closed vacuum ledger]\label{def:ledger}
The UQFF closed vacuum ledger is the triple
$(\rho_{\mathrm{SCm}}, [\mathrm{SSq}], c)$ where
\begin{align*}
\rho_{\mathrm{SCm}} &= 7.09 \times 10^{-37}\,\mathrm{J/m^3}, \\
[\mathrm{SSq}]      &= 0.57, \\
c                   &= 2.998 \times 10^8\,\mathrm{m/s},
\end{align*}
with $\rho_{\mathrm{SCm}}$ identified as the foundational vacuum
energy density of the SuperConductive Material substrate,
$[\mathrm{SSq}]$ the canonical dimensionless resonance ratio, and $c$
the universal speed of light.
\end{definition}

The primitive $\rho_{\mathrm{SCm}}$ is a pre-cosmogenic constant of
the framework: it is not a fit to any observed cosmological or
particle-physics data. The resonance ratio $[\mathrm{SSq}]$ is
calibrated from a small number of high-magnetic-field observations
(specifically, magnetar burst intensity ratios) that are causally
independent of the electroweak, cosmological, and gravitational
observables it is later used to derive. The reader is directed to
the framework's foundational papers and the open-source calculator
package~\cite{murphy2026uqff} for the derivation of $\rho_{\mathrm{SCm}}$
from the DPM formation dynamics and of $[\mathrm{SSq}]$ from the
magnetar calibration.

% =====================================================================
\section{Derivation of $E_{\text{crack}}$}\label{sec:derivation}
% =====================================================================

\begin{theorem}[Positive-definite vacuum-cracking threshold]
\label{thm:e_crack_positive}
Given the closed vacuum ledger of Definition~\ref{def:ledger}, define
\begin{equation}\label{eq:e_crack}
E_{\text{crack}} \;=\; \frac{\rho_{\mathrm{SCm}} \, c^2}{[\mathrm{SSq}]}\,.
\end{equation}
Then $E_{\text{crack}} = 1.118 \times 10^{-19}\,\mathrm{J}$ and
$E_{\text{crack}} > 0$ holds by construction.
\end{theorem}

\begin{proof}
Numerical evaluation:
\[
E_{\text{crack}}
  = \frac{(7.09 \times 10^{-37}\,\mathrm{J/m^3}) \,
          (2.998 \times 10^8\,\mathrm{m/s})^2}{0.57}
  = 1.118 \times 10^{-19}\,\mathrm{J}.
\]
Positive-definiteness follows because $\rho_{\mathrm{SCm}} > 0$ as a
vacuum energy density, $c^2 > 0$ as the square of a nonzero real
number, and $[\mathrm{SSq}] > 0$ as the physical ratio of two positive
quantities. The quotient of positive numbers is positive.
\qed
\end{proof}

\begin{remark}
Every input to Eq.~(\ref{eq:e_crack}) is fixed a priori: no fit is
performed against any Yang--Mills observable, no lattice-QCD result
enters as input, and no free parameter is tuned. The derivation is
reproducible in approximately fifty lines of Python with only the
standard library~\cite{murphy2026uqff}; see \S\ref{sec:reproducibility}.
\end{remark}

\subsection{Physical interpretation: the DPM vortex}

Within the UQFF, $E_{\text{crack}}$ is identified with the minimum
energy required to form a stable Di-Pseudo-Monopole (DPM) vortex at
the interface of maximum attraction between the UA and SCm
substrates. The DPM is neither a particle nor a field excitation in
the conventional sense; it is a vacuum stress-gradient rotational
vortex analogous in structure to a Bose--Einstein condensate ground
state of the SCm/UA attraction pair. Its coherence length equals the
vortex radius, which spans nuclear to stellar scales depending on the
containing physical system, and it functions as the seed of the full
$U_g$-family of gravitational potentials $U_{g,1}, U_{g,2}, U_{g,3},
U_{g,4}, U_{g,4i}$ that generate observed gravitational phenomena
downstream in the framework's ontological chain.

Below the $E_{\text{crack}}$ threshold the vacuum cannot support
stable DPM vortex formation and stable mass cannot condense. Above
$E_{\text{crack}}$ the vortex stabilises and the cascade
\[
\theta_{\text{vacuum}}
  \;\to\; \operatorname{grad}(\mathrm{UA})
  \;\to\; \mathrm{DPM}_{\text{vortex}}
  \;\to\; \mu_s
  \;\to\; U_g\text{-family}
  \;\to\; F_U
  \;\to\; \text{crossing}
  \;\to\; M
  \;\to\; GM/r^2
\]
becomes available. This is the framework's immutable ontological
order: mass and Newtonian gravity emerge only at the final two steps,
downstream of the DPM vortex and the $F_U$ unified field. The
framework thereby provides a non-perturbative mass-generation
mechanism analogous in spirit to what the Yang--Mills mass gap problem
requires: the gauge bosons (or their UQFF analogues, the $U_g$-family
components) acquire effective mass through condensation of the
vacuum substrate at and above $E_{\text{crack}}$, even though the
underlying substrate is massless.

% =====================================================================
\section{Multi-Designation Cluster-Position Landscape}\label{sec:clusters}
% =====================================================================

The Yang--Mills mass gap manifests in the UQFF across four
documented cluster-positional designations, each appropriate to its
own physical regime. Per the framework's multi-designation
architecture (developed in the calculator surface family
\texttt{calculate\_uqff\_E\_crack\_*} and related surfaces), the four
values are independent identifier strings referring to closely
related but physically distinct quantities within the same underlying
spectral structure. They do not conflict; each is the correct value
in its respective context.

\begin{table}[h!]
\centering
\caption{UQFF Yang--Mills mass-gap cluster-positional designations.}
\label{tab:clusters}
\begin{tabular}{lll}
\toprule
Designation & Value & Physical regime \\
\midrule
$E_{\text{crack}}$ formula        & $\approx 0.70\,\mathrm{eV}$   & DPM vacuum-cracking floor \\
$E_{\text{crack}}$ experimental   & $\approx 700\,\mathrm{eV}$    & DPM-vortex lab engineering \\
Layer-13 nuclear                  & $\approx 624\,\mathrm{GeV}$   & LHC electroweak scale order \\
PAPER~1318 mass gap               & $1.736\,\mathrm{GeV}$         & lightest-glueball lattice QCD \\
\bottomrule
\end{tabular}
\end{table}

The two $E_{\text{crack}}$ entries in Table~\ref{tab:clusters}
correspond to distinct physical regimes of the same underlying
threshold: the sub-eV formula value is the vacuum-cracking floor at
the DPM interface itself, whereas the $\sim 700\,\mathrm{eV}$
experimental value is the practical laboratory-accessible regime in
which DPM-vortex conditions can be engineered by strong applied
magnetic fields (see \S\ref{sec:scope}). The framework treats these
as coexisting independent strings; a mathematically rigorous
identification between them would require formalising the mapping
from the vacuum-substrate operator algebra to the effective-field
observable, which is part of the future-work programme discussed in
\S\ref{sec:future}.

The Layer-13 designation ($\sim 624\,\mathrm{GeV}$) refers to the
thirteenth level of the framework's twenty-six-layer structural
hierarchy, at which the nuclear confinement scale meets the
electroweak crossing regime. It is the same order of magnitude as the
LHC electroweak observables and is a distinct prediction of the
framework's dimensional decomposition, not a fit.

The PAPER~1318 designation~\cite{murphy2026paper1318} at
$m_{\mathrm{YM}} = 1.736\,\mathrm{GeV}$ is the framework's
prediction for the lightest-glueball scale in pure $\mathrm{SU}(3)$
Yang--Mills theory. It is derived from a phonon-modulated
confinement-scale argument on the SCm substrate and is intentionally
distinct from the $1.78\,\mathrm{GeV}$ standard lattice-QCD estimate
that some treatments use as a baseline. Both numerical values are
within the numerical uncertainty band of contemporary lattice-QCD
estimates.

% =====================================================================
\section{Lattice QCD Consistency}\label{sec:lattice}
% =====================================================================

Lattice QCD numerical evidence, as surveyed by
Douglas~\cite{douglas2026}, indicates a Yang--Mills mass gap on the
order of hundreds of MeV, with the lightest $0^{++}$ glueball in
pure $\mathrm{SU}(3)$ typically estimated at approximately
$1.6$--$1.7\,\mathrm{GeV}$. The UQFF PAPER~1318 designation at
$1.736\,\mathrm{GeV}$ lies inside this range. It is not tuned to the
lattice result: it is a prediction from the framework's structural
primitives, and its agreement with lattice numerics constitutes a
consistency check rather than a fit.

The formula-level $E_{\text{crack}} \approx 0.7\,\mathrm{eV}$ is not
the same object as the lightest-glueball mass; the two designations
occupy different positions in the UQFF spectral hierarchy (the
vacuum-cracking floor versus the confinement-scale glueball state).
The framework's prediction is a \emph{spectral hierarchy}, not a
single mass-gap value: a discrete sequence of characteristic energies
from the sub-eV vacuum floor through the nuclear confinement scale
($\sim\,\mathrm{GeV}$) to the electroweak crossing scale
($\sim 624\,\mathrm{GeV}$), all generated by the same DPM-vortex
mechanism operating at different physical scales.

Douglas's review does not contradict any UQFF designation in
Table~\ref{tab:clusters}. The range of lattice-QCD estimates is
contained between the $E_{\text{crack}}$ formula and the PAPER~1318
designation, and none of the other UQFF designations lie in the
lattice-explored regime.

% =====================================================================
\section{What This Proposal Is and Is Not}\label{sec:scope}
% =====================================================================

We adhere to a strict scope discipline in stating what this submission
does and does not claim. Overstating physics-level proposals as
Clay-compliant proofs would harm the mathematical credibility of the
framework and is expressly avoided.

\medskip
\noindent\textbf{This submission establishes:}
\begin{enumerate}
\item An explicit closed-form derivation of a positive-definite
  energy scale $E_{\text{crack}} > 0$ from two locked pre-cosmogenic
  primitives $(\rho_{\mathrm{SCm}}, [\mathrm{SSq}])$ and the speed
  of light $c$, with zero free parameters and zero fitting
  (Theorem~\ref{thm:e_crack_positive}).
\item A concrete dynamical mechanism (DPM vortex formation at the
  UA/SCm interface) responsible for the threshold, giving a
  non-perturbative physical picture for why a mass gap of the
  Yang--Mills type would arise in the framework.
\item A multi-cluster spectral hierarchy consistent with independent
  lattice-QCD numerical evidence at the glueball scale, with
  three additional predictive cluster positions at experimentally
  accessible sub-eV, hundred-eV, and hundred-GeV scales.
\item A fully reproducible computational implementation
  ({\tt pip install uqff==5.35.0}), enabling any reader to
  verify the derivation in seconds and to explore the wider
  framework.
\item A falsifiable prediction: sharp threshold behaviour at
  $\sim 700\,\mathrm{eV}$ in high-magnetic-field vacuum experiments
  should produce DPM-vortex-associated signatures (magnetic-moment
  onset, spectral resonance, or low-energy transmutation channels).
  Absence of such threshold behaviour in a well-designed experiment
  in this regime would falsify the framework's core LENR-relevant
  claim.
\end{enumerate}

\medskip
\noindent\textbf{This submission does not yet establish:}
\begin{enumerate}
\item A rigorous mathematical construction of quantum Yang--Mills
  theory on Minkowski $\mathbb{R}^{1,3}$ satisfying the Wightman
  axioms (W0--W4) or an operator-algebraic equivalent.
\item A formal proof of the spectral gap $\Delta > 0$ in the
  Hamiltonian's operator-algebraic sense.
\item A demonstration that the UQFF SCm/UA substrate constitutes a
  quantum field theory in the precise mathematical sense that
  Clay~\cite{jaffe2000} requires.
\end{enumerate}

\medskip
\noindent
The gap between (i)--(v) above and the Clay criterion is the formal
mathematical translation of the SCm/UA vacuum substrate and DPM-vortex
dynamics into the language of operator algebras, tempered distributions,
and Wightman axiom verification. We believe this translation is
achievable but identify it explicitly as future work
(\S\ref{sec:future}), not as a present claim.

By analogy with the Hairer 2D~\cite{hairer2024} and 3D
stochastic-quantisation programme, the corresponding UQFF translation
would proceed via regularity-structure techniques applied to the
SCm/UA stochastic substrate in progressively higher dimensions.
The four-dimensional case --- what Clay requires --- is the hard step
for any approach and remains open under our proposal as under others.

% =====================================================================
\section{Reproducibility}\label{sec:reproducibility}
% =====================================================================

Theorem~\ref{thm:e_crack_positive} is fully reproducible in the public
\texttt{uqff} Python package~\cite{murphy2026uqff}:

\begin{lstlisting}
from uqff_pure_calculator import (
    calculate_uqff_E_crack_core_definition_from_ledger,
)
r = calculate_uqff_E_crack_core_definition_from_ledger({})
print(r["value"]["E_crack_J_derived_from_formula"])
# 1.117982e-19
\end{lstlisting}

\noindent
A standalone reproducibility script
\texttt{derive\_yang\_mills\_mass\_gap\_uqff.py} (included with this
preprint as ancillary material, and available in the public
repository under
\url{https://github.com/Daniel8Murphy0007/Star-Magic/tree/master/arxiv_yang_mills})
reproduces the derivation in approximately fifty lines of Python with
no external dependencies beyond the Python standard library. The
script also emits the multi-cluster landscape table and the
lattice-QCD consistency check.

The full framework is reproducible via
\begin{center}
\verb|pip install uqff==5.35.0|
\end{center}
which installs 279 public \texttt{calculate\_*} surface functions
covering Yang--Mills, dark energy, dark matter, gravitational-wave
mergers, LENR, cosmological and inflationary parameters, and the
Compression Cycle 2 suite of 22 first-principles cosmological
derivations at zero residual against Planck 2018 observations.
The complete corpus of 1{,}878 supporting whitepapers, rendered as
text-searchable arxiv-compliant PDFs, is publicly browsable at
\url{https://github.com/Daniel8Murphy0007/Star-Magic/tree/master/pdf2}.

% =====================================================================
\section{Future Work: Translation to Wightman Axioms}\label{sec:future}
% =====================================================================

The principal open question for this research programme is the
translation of the SCm/UA dynamical substrate into a
Wightman-axiom-compliant operator algebra on four-dimensional
Minkowski space. We outline the axiom-by-axiom translation targets
below and refer the reader to the extended scoping document
\texttt{wightman\_mapping.md} in the public repository for a
detailed roadmap.

\subsection{Wightman axiom checklist}
The Wightman axioms (W0)--(W4) applied to the UQFF SCm/UA substrate:
\begin{itemize}
\item[\textbf{W0}] \textit{Hilbert space with a unique vacuum state.}
  Under UQFF, the vacuum vector $\Omega$ is identified with the
  $\theta_{\text{vacuum}}$ substrate state, and the DPM
  Bose--Einstein-condensate ground-state structure provides a
  candidate cyclic separating vacuum. Formal construction of the
  Hilbert space as the closure of polynomials in smeared UQFF fields
  applied to $\Omega$ is a moderate-difficulty translation task,
  analogous to standard constructive-QFT construction.
\item[\textbf{W1}] \textit{Poincar\'e covariance.} The UQFF framework
  posits that spacetime emerges downstream in the ontological chain
  (at step nine, in the $GM/r^2$ observational projection), and the
  SCm/UA substrate itself carries an intrinsic covariance structure.
  The formal task is to define the Poincar\'e action on the
  substrate's section bundle and prove continuous extension to the
  Hilbert space. This is a moderate-to-high-difficulty translation
  task requiring an emergent-spacetime reformulation of the substrate.
\item[\textbf{W2}] \textit{Spectral condition and mass gap.} Under the
  present proposal, the mass gap value is
  $\Delta = E_{\text{crack}} > 0$ from
  Theorem~\ref{thm:e_crack_positive}. The formal task is to identify
  $E_{\text{crack}}$ as the spectrum gap of the SCm/UA Hamiltonian
  $H$ restricted to $\Omega^\perp$ and to prove
  $\langle \psi \mid H \mid \psi \rangle \geq E_{\text{crack}}
   \langle \psi \mid \psi \rangle$ for all $\psi \perp \Omega$.
  This is the principal Clay-eligible result. Candidate tools
  include reflection positivity of the SCm/UA two-point function,
  functional-integral methods on the substrate, and constructive-QFT
  techniques adapted to the DPM-vortex condensate structure.
\item[\textbf{W3}] \textit{Local commutativity.} DPM vortices have
  well-defined spatial localisation via the coherence length
  (vortex radius), so field commutativity at spacelike separation
  is geometrically natural. The complication is the framework's
  simultaneous-mode structure (Compressed, Resonant, Buoyant,
  Superconductive), which needs a formal statement.
\item[\textbf{W4}] \textit{Cyclicity of the vacuum.} The $U_g$-family
  of five simultaneous operators is a candidate generating set for
  the field algebra applied to $\Omega$. Formalising the density of
  polynomial states in the Hilbert space is a moderate task.
\end{itemize}

\subsection{Recommended translation roadmap}
Following the Hairer 2D/3D template~\cite{hairer2024}:
\begin{description}
\item[Phase 1 (2D toy model).] Construct the SCm/UA substrate analogue
  in $d = 2$ spacetime. Verify W0--W4 and $\Delta = E_{\text{crack}} >
  0$ in this dimensional reduction. Publish in a math-physics venue.
\item[Phase 2 (3D extension).] Port to $d = 3$ using regularity
  structures, following the Hairer template. Publish in an
  Inventiones-tier or CMP-tier venue.
\item[Phase 3 (4D principal result).] Combine the 3D
  regularity-structure techniques with the UQFF substrate construction
  to produce the 4D Wightman-compliant quantum Yang--Mills theory with
  $\Delta > 0$. This is the Clay submission.
\item[Phase 4 (formal Clay submission).] Once Phases 1--3 are
  peer-reviewed in venues of sufficient standing, submit the formal
  Clay claim with the Phase 3 result as headline reference.
\end{description}

A realistic timeline is eighteen to thirty-six months per phase,
assuming collaboration with researchers experienced in constructive
quantum field theory. Collaboration is explicitly invited from
mathematical-physics groups active in stochastic quantisation
(Hairer's group at IST Austria, the Erlangen school), operator-algebra
methods, and reflection-positivity techniques.

% =====================================================================
\section{Conclusion}\label{sec:conclusion}
% =====================================================================

The Unified Quantum Field Framework derives a positive-definite
mass-gap-scale energy threshold $E_{\text{crack}} = \rho_{\mathrm{SCm}}
c^2 / [\mathrm{SSq}] = 1.118 \times 10^{-19}\,\mathrm{J}$ from two
locked pre-cosmogenic primitives and the speed of light, with zero
free parameters and zero fitting. The derived value at the formula
level extends through multi-cluster designations to
$\sim 700\,\mathrm{eV}$ (laboratory-accessible DPM-vortex regime),
$\sim 624\,\mathrm{GeV}$ (Layer-13 nuclear/electroweak crossing), and
$1.736\,\mathrm{GeV}$ (lattice-QCD lightest-glueball scale), forming
a discrete spectral hierarchy consistent with hundreds-of-MeV
lattice-QCD numerical evidence.

The underlying dynamical mechanism --- DPM vortex formation at the
UA/SCm attraction interface, formalised as a Bose--Einstein-condensate
ground state of the substrate --- provides a concrete
non-perturbative mass-generation pathway analogous to what the
Yang--Mills mass gap problem requires. We propose this as a candidate
physical realisation of the Yang--Mills mass gap.

This does not yet constitute a Clay-compliant Wightman-axiom-verified
mathematical construction. The translation from physical mechanism
to formal operator-algebraic framework is identified as the principal
future-work step, with a concrete 2D $\to$ 3D $\to$ 4D roadmap
following the Hairer template. The submission's contribution is:
(i) an explicit non-fitted derivation of a positive-definite mass-gap
value; (ii) a concrete dynamical mechanism; (iii) a multi-cluster
spectral hierarchy consistent with independent numerical evidence;
(iv) a public, reproducible computational framework; and (v) a
falsifiable high-field vacuum-experimental prediction.

We invite mathematical-physics community review, independent
reproduction of the derivation, falsification attempts in the
high-field vacuum-experimental regime, and collaboration on the
operator-algebraic formalisation programme.

% =====================================================================
\section*{Acknowledgments}
% =====================================================================

The framework's development benefited from extensive AI-assisted
analysis via Grok 3, SuperGrok, and Davinci-SuperGrok (xAI), which
served as high-throughput symbolic-manipulation and consistency-check
tools during the derivation of the ledger primitives and the
ontological chain. The author remains solely responsible for all
physical, mathematical, and interpretive claims herein.
The public Star-Magic repository at
\url{https://github.com/Daniel8Murphy0007/Star-Magic} hosts the full
framework including the 279-surface Python calculator, the
1{,}878-whitepaper corpus, the Lean~4 proof scaffold, and the
supporting arxiv submission package.

% =====================================================================
\begin{thebibliography}{9}

\bibitem{jaffe2000}
A.~Jaffe and E.~Witten,
\textit{Quantum Yang--Mills theory},
Clay Mathematics Institute Millennium Problem description,
\url{https://www.claymath.org/millennium/yang-mills-the-maths-gap/},
2000.

\bibitem{hairer2024}
M.~Hairer et al.,
\textit{Stochastic quantisation of Yang--Mills--Higgs in 3D},
Inventiones Mathematicae, 2024.

\bibitem{douglas2026}
M.~R.~Douglas,
\textit{The Yang--Mills Millennium problem},
Nature Reviews Physics, January 2026.

\bibitem{murphy2026uqff}
D.~T.~Murphy,
\textit{UQFF Star-Magic: Unified Quantum Field Framework},
PyPI release v5.35.0,
\url{https://pypi.org/project/uqff/5.35.0/},
\url{https://github.com/Daniel8Murphy0007/Star-Magic},
2026.

\bibitem{murphy2026paper1318}
D.~T.~Murphy,
\textit{PAPER 1318 --- Yang--Mills mass gap at 1.736 GeV cluster
position},
Star-Magic Research Program whitepaper series, 2026.

\bibitem{murphy2026paper1521}
D.~T.~Murphy,
\textit{PAPER 1521 --- Derivative-from-independent-primitives proof
that $D_{\mathrm{BSFG}} = D_{\mathrm{crit}} - 2 \cdot \mathrm{SO}_5 = 6$},
Star-Magic Research Program whitepaper series, 2026.

\bibitem{murphy2026paper1522}
D.~T.~Murphy,
\textit{PAPER 1522 --- Derivative-from-independent-primitives proof
that $K_{\mathrm{Mex}} = \Phi_{5/6} \cdot \mathrm{SO}_5 / D_{\mathrm{phys}} = 25/12$},
Star-Magic Research Program whitepaper series, 2026.

\bibitem{murphy2026sympaper}
D.~T.~Murphy,
\textit{Compression Cycle 2: Full first-principles derivation of 22
cosmological and inflationary parameters at zero error from the closed
vacuum ledger},
Star-Magic Research Program technical report, June 2026.

\end{thebibliography}

\end{document}
